\documentclass[11pt]{article}

\usepackage[T1]{fontenc}
\usepackage[utf8]{inputenc}
\usepackage{amsmath,amssymb,amsthm}
\usepackage{array}
\usepackage{booktabs}
\usepackage{geometry}
\usepackage{hyperref}
\usepackage{longtable}
\usepackage{microtype}
\usepackage{xcolor}

\geometry{margin=0.92in}
\hypersetup{
  colorlinks=true,
  linkcolor=blue!55!black,
  citecolor=blue!55!black,
  urlcolor=blue!55!black
}

\newtheorem{theorem}{Theorem}[section]
\newtheorem{lemma}[theorem]{Lemma}
\newtheorem{corollary}[theorem]{Corollary}
\newtheorem{proposition}[theorem]{Proposition}
\theoremstyle{definition}
\newtheorem{definition}[theorem]{Definition}
\theoremstyle{remark}
\newtheorem{remark}[theorem]{Remark}

\newcommand{\F}{\mathbb F}
\newcommand{\xor}{\mathbin{\oplus}}
\newcommand{\FiveCDC}{\textup{FiveCDC}}
\newcommand{\Dfive}{D_5}

\title{Two Exact Four-Pole Reductions for the\\
Standard Five-Cycle Double Cover Conjecture}
\author{Atharva Vaidya\\
\small AI-assisted working preprint for independent human verification}
\date{July 28, 2026}

\begin{document}
\maketitle

\begin{center}
\fcolorbox{red!65!black}{red!4}{
\parbox{0.91\textwidth}{
\textbf{Resolution status.}
The standard five-cycle double cover conjecture remains open.
This note proves two positive substitution reductions.  It proves neither
the conjecture nor its negation, and it makes no claim about the orientable
five-cycle double cover conjecture.
}}
\end{center}

\begin{center}
\fcolorbox{orange!70!black}{yellow!10}{
\parbox{0.91\textwidth}{
\textbf{AI-use disclosure.}
OpenAI Codex agents, under Atharva Vaidya's direction, selected and tested
the poles, found the finite label certificates, wrote the checking programs,
and drafted substantial portions of this manuscript.  The certificate
tables and source code are supplied for independent human checking.
AI assistance is not peer review.  No human specialist has yet certified
the literature comparison or accepted scholarly responsibility for a
submission.
}}
\end{center}

\begin{abstract}
Write \(\Dfive=\binom{[5]}2\subseteq\F_2^5\).  A standard five-cycle
double cover of a cubic graph is equivalently a \(\Dfive\)-label on every
edge whose three incident labels xor to zero at each vertex.  This makes a
cubic four-pole an exact finite boundary relation.

We give two human-checkable positive relations.  First, the four-pole
obtained by deleting adjacent vertices from the Petersen graph extends
every boundary word that is a permutation of \((q,q,r,r)\), for arbitrary
\(q,r\in\Dfive\).  A 13-row certificate covers the three intersection
types of \(q,r\) and all port placements.  Second, a four-pole obtained by
deleting adjacent vertices from the Heawood graph has the full boundary
relation: it extends every ordered
\((A_0,A_1,A_2,A_3)\in\Dfive^4\) with xor zero.  A 10-row certificate,
together with the global \(S_5\) action, covers all 640 admissible words.

It follows that replacing two independent edges of any FiveCDC-positive
cubic graph by either pole preserves FiveCDC for every bijection from the
four exposed endpoints to the four ports.  We prove separately that the
operation preserves simplicity and bridgelessness under the stated
hypotheses.  Thus neither substitution can manufacture a minimum
counterexample from a smaller positive graph.  The Petersen certificate
establishes only the repeated-pair interface needed by this insertion,
whereas the Heawood certificate proves the complete xor-zero boundary
relation.

The tempting extension to all bridgeless bipartite cubic graphs is false:
deleting adjacent vertices from \(K_{3,3}\) gives \(K_{2,2}\), whose
ordered boundary relation has 580 of the 640 admissible words and misses
one 60-word orbit.  We give a direct proof of that obstruction.  All claims
here are local positive reductions or scope boundaries, not a resolution
of FiveCDC.
\end{abstract}

\tableofcontents

\section{Status, conventions, and prior-art boundary}

A standard five-cycle double cover of a finite graph is a family of at
most five Eulerian edge-subsets in which every edge occurs exactly twice.
Equivalently, one may require exactly five indexed subsets and allow empty
members.  An Eulerian edge-subset need not be connected or 2-regular.
Recent papers continue to state FiveCDC as open
\cite{Oum2026,HusekSamal2026}; special cases are known
\cite{LiuEtAl2023}.

Graphs in the two insertion theorems below are finite and undirected.
The input and output graphs are simple and cubic.  The pair-label
equivalence itself also works for multigraphs if parallel edges are kept
as distinct edge objects and a loop contributes two incidences at its
endpoint.  The finite pole certificates are for the displayed simple
cores only; no computational conclusion is extrapolated to arbitrary
multigraphs.

M\'a\v cajov\'a, Mazzuoccolo, and Tabarelli introduced the relevant
CDC-colouring language for multipoles and the ten boundary types of a
four-pole \cite{MacajovaEtAl2026}.  That boundary framework, including its
use in minimum-counterexample reductions, is prior work and is not claimed
here.  The present contribution is the two explicit extension certificates
and the resulting insertion lemmas.  We have not completed a specialist
search capable of deciding whether either particular pole theorem, or a
more general theorem implying it, already occurs in the literature.
Accordingly no novelty or priority claim is made.  The authors of the
four-pole literature should be consulted before public submission.

The orientable five-cycle double cover conjecture adds a direction
condition: every edge must be traversed once in each direction.  No
orientation variables or directed compatibility conditions occur below,
so none of the results imply an orientable statement.

\section{Pair labels and exact boundary relations}

Put
\[
  [5]=\{0,1,2,3,4\},\qquad
  \Dfive=\binom{[5]}2\subseteq\F_2^5.
\]
We identify a pair with its weight-two characteristic vector.

\begin{proposition}[pair-label equivalence]\label{prop:labels}
For a finite graph \(G\), the following are equivalent.
\begin{enumerate}
\item There are five Eulerian edge-subsets
      \(C_0,\ldots,C_4\) containing every edge exactly twice.
\item Every edge has a label \(L(e)\in\Dfive\), and the xor of the labels
      on the edge incidences at every vertex is zero.
\end{enumerate}
\end{proposition}

\begin{proof}
Given the five subsets, set
\[
                   L(e)=\{i:e\in C_i\}.
\]
Exact double coverage gives \(|L(e)|=2\).  In coordinate \(i\), the
vertex xor equation is exactly the even-degree condition for \(C_i\).
Conversely, the coordinate supports of a pair labelling are Eulerian by
the vertex equations, and every edge belongs to its two label coordinates.
\end{proof}

At a cubic vertex, three labels \(A,B,C\in\Dfive\) have
\(A\xor B\xor C=0\) if and only if they are the three edges of a triangle
in \(K_5\).

\begin{definition}
A \emph{terminal-distinct cubic four-pole} is a proper graph with four
specified degree-two vertices, called ports, and every other vertex of
degree three.  Completing each port with one semiedge makes every vertex
cubic.  Its ordered boundary relation \(\mathcal R_5(P)\) is the set of
four semiedge-label words that extend to a \(\Dfive\)-labelling of all
proper edges with xor zero at every completed vertex.
\end{definition}

\begin{lemma}[boundary parity]\label{lem:boundary-parity}
Every \((A_0,A_1,A_2,A_3)\in\mathcal R_5(P)\) satisfies
\[
                         A_0\xor A_1\xor A_2\xor A_3=0.       \tag{2.1}
\]
\end{lemma}

\begin{proof}
Xor all internal vertex equations.  Each proper edge occurs at two
endpoints and cancels.  Only the four semiedge labels remain.
\end{proof}

\begin{lemma}[ten full-boundary types]\label{lem:ten-types}
The 640 ordered words in \(\Dfive^4\) satisfying \((2.1)\) form ten orbits
under the simultaneous \(S_5\) action on coordinates.
\end{lemma}

\begin{proof}
View the four pair labels as four ordered edges of a loopless multigraph on
the five coordinate vertices.  Equation~(2.1) says that every coordinate
has even degree.  Ignoring edge order, the possibilities are four parallel
edges, two doubled adjacent pairs, two doubled disjoint pairs, or one
simple four-cycle.  The first has one ordered-edge orbit.  Each of the
other three shapes has three orbits, determined by the induced partition
or opposite-position pair among the four boundary positions.  Hence there
are \(1+3+3+3=10\) orbits.  Their sizes in the table of
Section~\ref{sec:heawood} sum to 640.
\end{proof}

\section{A generic insertion lemma}

Let \(P\) be a connected terminal-distinct cubic four-pole with ordered
port set \(T\).  Let \(G\) be a cubic graph and let
\[
                         e=ab,\qquad f=cd
\]
be independent edges.  For an arbitrary bijection
\(\pi:\{a,b,c,d\}\to T\), define
\[
        G\star_{\pi}P
        =\bigl(G-\{e,f\}\bigr)\cup P
          \cup\{z\pi(z):z\in\{a,b,c,d\}\}.              \tag{3.1}
\]
There is no restriction on which ports receive the two ends of one old
edge.

\begin{definition}
The pole \(P\) is \emph{repeated-pair universal} if, for all
\(q,r\in\Dfive\), every ordered placement of the multiset
\((q,q,r,r)\) belongs to \(\mathcal R_5(P)\).  It is
\emph{full-boundary} if equality holds in Lemma~\ref{lem:boundary-parity}.
\end{definition}

Every full-boundary pole is repeated-pair universal.

\begin{theorem}[positive insertion]\label{thm:insertion}
If \(G\) has a standard FiveCDC and \(P\) is repeated-pair universal, then
\(G\star_\pi P\) has a standard FiveCDC for every port bijection \(\pi\).
\end{theorem}

\begin{proof}
Take a \(\Dfive\)-labelling of \(G\), and let \(q=L(e)\) and \(r=L(f)\).
Label the new joining edges at \(a,b\) by \(q\) and those at \(c,d\) by
\(r\).  In the ordered port convention this is some placement of
\((q,q,r,r)\), so it extends through \(P\).  Every old vertex retains
the same local equation it had in \(G\).  Proposition~\ref{prop:labels}
converts the resulting labelling into a standard FiveCDC.
\end{proof}

\begin{lemma}[simplicity and bridgelessness]\label{lem:bridgeless}
Suppose \(G\) is connected, simple, cubic, and bridgeless; \(e,f\) are
independent; and the proper core of \(P\) is connected, simple, and
bridgeless.  Then \(G\star_\pi P\) is connected, simple, cubic, and
bridgeless for every \(\pi\).
\end{lemma}

\begin{proof}
Cubicity and simplicity follow directly from fresh, terminal-distinct pole
vertices and four distinct exposed endpoints.

Put \(H=G-\{e,f\}\).  Every component of \(H\) contains at least two
exposed endpoints: a component containing exactly one would have a
one-edge boundary in \(G\), making that edge a bridge.  Joining all
components to the connected pole makes the output connected.

Every proper pole edge lies on a circuit inside the bridgeless core.  For a
joining edge \(z\pi(z)\), choose another exposed endpoint \(z'\) in the
same component of \(H\).  A \(z\)-to-\(z'\) path in that component, the
joining edge at \(z'\), and a path in the pole from \(\pi(z')\) to
\(\pi(z)\) close a circuit through \(z\pi(z)\).

Finally let \(g\) be an old edge of \(H\).  If it is not a bridge in its
\(H\)-component, it already lies on a circuit.  Otherwise deleting \(g\)
splits that component into shores \(X,Y\).  Each shore contains an exposed
endpoint; if \(X\) did not, then \(\delta_G(X)=\{g\}\), contradicting
bridgelessness of \(G\).  Paths in \(X,Y\), two joining edges, and a path
through the pole close a circuit through \(g\).  Every output edge is
therefore on a circuit.
\end{proof}

\begin{corollary}[minimal-counterexample pruning]\label{cor:prune}
If a graph made by the insertion \((3.1)\) is a FiveCDC counterexample,
then the smaller graph \(G\) is already a counterexample.  In particular,
neither pole certified below can create a minimum counterexample from a
smaller FiveCDC-positive graph.
\end{corollary}

\section{The Petersen repeated-pair certificate}\label{sec:petersen}

Use the Petersen graph6 record
\[
                         \texttt{ICOf@pSb?}
\]
and delete the adjacent vertices \(0,3\).  Keep the original vertex names.
The ordered ports are
\[
                         (6,9,7,8),                       \tag{4.1}
\]
and the ordered proper edges are
\[
\begin{aligned}
 &(1,4),(1,6),(1,8),(2,5),(2,6),\\
 &(2,7),(4,7),(4,9),(5,8),(5,9).
\end{aligned}                                             \tag{4.2}
\]
The proper core is simple, connected, bridgeless, and has girth five.

\begin{theorem}[Petersen repeated-pair extension]\label{thm:petersen}
For all \(q,r\in\Dfive\), every ordered placement of
\((q,q,r,r)\) on the ports \((4.1)\) extends through the Petersen pole.
\end{theorem}

\begin{proof}
Under \(S_5\), an ordered pair \(q,r\) is of one of three kinds: equal,
distinct with one common coordinate, or disjoint.  Representatives are
\((01,01)\), \((01,02)\), and \((01,23)\).  For unequal labels there are
six placements of the two copies of each label.  The following 13 rows
give the boundary labels in port order \((4.1)\), followed by labels on
the ten proper edges in order \((4.2)\).

\begingroup
\setlength{\tabcolsep}{4pt}
\renewcommand{\arraystretch}{1.08}
\small
\begin{longtable}{@{}>{\ttfamily\raggedright\arraybackslash}p{0.22\textwidth}
                       >{\ttfamily\raggedright\arraybackslash}p{0.70\textwidth}@{}}
\toprule
\normalfont Port labels & \normalfont Proper-edge labels\\
\midrule
\endfirsthead
\toprule
\normalfont Port labels & \normalfont Proper-edge labels\\
\midrule
\endhead
01 01 01 01 & 01 02 12 01 12 02 12 02 02 12\\
01 01 02 02 & 01 02 12 02 12 01 12 02 01 12\\
01 02 01 02 & 02 03 23 34 13 14 04 24 03 04\\
01 02 02 01 & 34 03 04 01 13 03 23 24 14 04\\
02 01 01 02 & 34 03 04 02 23 03 13 14 24 04\\
02 01 02 01 & 01 03 13 34 23 24 04 14 03 04\\
02 02 01 01 & 02 01 12 01 12 02 12 01 02 12\\
01 01 23 23 & 01 02 12 01 12 02 03 13 13 03\\
01 23 01 23 & 01 02 12 01 12 02 12 02 13 03\\
01 23 23 01 & 01 02 12 01 12 02 03 13 02 12\\
23 01 01 23 & 01 02 12 01 03 13 03 13 13 03\\
23 01 23 01 & 01 02 12 01 03 13 12 02 02 12\\
23 23 01 01 & 01 02 12 01 03 13 03 13 02 12\\
\bottomrule
\end{longtable}
\endgroup

Every entry has weight two.  At each of the eight completed pole vertices,
the xor of its three incident labels is zero; this is checked directly
from \((4.1)\)--\((4.2)\).  Applying a simultaneous coordinate
permutation preserves these equations.  The three intersection types and
the displayed placements exhaust all \(q,r\), proving the theorem.
\end{proof}

\begin{remark}[certified scope]
The 13-row theorem proves exactly the interface required by
Theorem~\ref{thm:insertion}.  It is not presented here as a classification
of all 640 xor-zero boundary words.  The Heawood theorem below has that
stronger, full-boundary conclusion.
\end{remark}

\section{The Heawood full-boundary certificate}\label{sec:heawood}

The Heawood four-pole has graph6 record
\[
                         \texttt{KhEGHC@AI?\_P}.
\]
Its ordered ports are
\[
                         (0,3,8,11),                       \tag{5.1}
\]
and its ordered proper edges are
\[
\begin{aligned}
 &(0,1),(1,2),(2,3),(3,4),(0,5),(4,5),(5,6),(2,7),\\
 &(6,7),(7,8),(4,9),(8,9),(1,10),(9,10),(6,11),(10,11).
\end{aligned}                                             \tag{5.2}
\]
This proper core is simple, connected, bridgeless, and has girth six.

\begin{theorem}[Heawood full boundary]\label{thm:heawood}
For every \(A_0,A_1,A_2,A_3\in\Dfive\) satisfying
\[
                         A_0\xor A_1\xor A_2\xor A_3=0,
\]
the Heawood four-pole has a \(\Dfive\)-labelling with those labels on the
ordered ports \((5.1)\).
\end{theorem}

\begin{proof}
The table contains one representative of each orbit in
Lemma~\ref{lem:ten-types}.  The middle column gives labels on the ports
\((5.1)\); the last column gives labels on the sixteen edges \((5.2)\).

\begingroup
\setlength{\tabcolsep}{3.5pt}
\renewcommand{\arraystretch}{1.08}
\scriptsize
\begin{longtable}{@{}l
 >{\ttfamily\raggedright\arraybackslash}p{0.18\textwidth}
 >{\ttfamily\raggedright\arraybackslash}p{0.66\textwidth}@{}}
\toprule
Type & \normalfont Port labels & \normalfont Proper-edge labels\\
\midrule
\endfirsthead
\toprule
Type & \normalfont Port labels & \normalfont Proper-edge labels\\
\midrule
\endhead
\texttt{AA} & 02 02 02 02 &
01 02 01 12 12 02 01 12 02 01 01 12 12 02 12 01\\
\texttt{AT2} & 02 02 03 03 &
01 02 03 23 12 02 01 23 03 02 03 23 12 02 13 01\\
\texttt{AT3} & 02 03 02 03 &
01 02 04 34 12 14 24 24 34 23 13 03 12 01 23 02\\
\texttt{AT4} & 02 03 03 02 &
01 02 01 13 12 01 02 12 01 02 03 23 12 02 12 01\\
\texttt{T2T2} & 02 02 13 13 &
01 02 01 12 12 02 01 12 02 01 01 03 12 13 12 23\\
\texttt{T2T3} & 02 03 12 13 &
01 02 01 13 12 01 02 12 01 02 03 01 12 13 12 23\\
\texttt{T2T4} & 02 03 13 12 &
01 02 04 34 12 13 23 24 12 14 14 34 12 13 13 23\\
\texttt{T3T3} & 02 13 02 13 &
01 02 01 03 12 02 01 12 02 01 23 12 12 13 12 23\\
\texttt{T3T4} & 02 13 03 12 &
01 02 03 01 12 02 01 23 03 02 12 23 12 13 13 23\\
\texttt{T4T4} & 02 13 13 02 &
01 02 01 03 12 01 02 12 24 14 13 34 12 14 04 24\\
\bottomrule
\end{longtable}
\endgroup

Every displayed label has weight two.  Completing the ports in order,
the xor of the three incident labels is zero at each of the twelve
vertices.  Thus every row is a direct positive certificate.

The orbit sizes in the displayed order are
\[
                 10,60,60,60,30,120,120,30,120,30,
\]
whose sum is 640.  Lemma~\ref{lem:ten-types} proves that these are all
admissible boundary orbits.  A simultaneous coordinate permutation of a
row preserves label weight and every vertex xor equation, so all 640 words
extend.
\end{proof}

\begin{corollary}[two non-covering reductions]\label{cor:two-poles}
The insertion theorem and the bridgelessness lemma apply to both the
Petersen and Heawood poles, with every one of the \(4!\) port bijections.
Neither construction is a graph cover or lift.
\end{corollary}

\section{A sharp scope boundary: \(K_{3,3}\)}

The Heawood graph is cubic, bipartite, and bridgeless.  It is therefore
natural to ask whether deleting the ends of any edge from any bridgeless
bipartite cubic graph always gives a full-boundary four-pole.  The smallest
example shows that this statement is false.

\begin{proposition}[\(K_{3,3}\) obstruction]\label{prop:k33}
Delete adjacent vertices from \(K_{3,3}\).  The proper four-pole is
\(K_{2,2}\).  Order its ports with the two vertices of one bipartition
first and the two vertices of the other bipartition second.  Its boundary
relation contains 580 of the 640 xor-zero words and misses exactly the
60-word orbit represented by
\[
                              (02,02,03,03).              \tag{6.1}
\]
In particular it is not full-boundary.
\end{proposition}

\begin{proof}
The exact count and uniqueness of the missing orbit follow from the
standard-library finite replay described in Section~\ref{sec:replay}.
Here is a direct proof that \((6.1)\) does not extend.

Let the two first-part vertices be \(x_0,x_1\), and let \(y\) be either
second-part vertex.  At \(x_0\), the two proper-edge labels have xor \(02\).
Therefore they are
\[
                         \{0,t_0\},\ \{2,t_0\}
\]
in some order, where \(t_0\notin\{0,2\}\).  The same reasoning at \(x_1\)
shows that its two proper-edge labels are
\[
                         \{0,t_1\},\ \{2,t_1\}
\]
in some order, with \(t_1\notin\{0,2\}\).

The two proper edges incident with \(y\), one from each first-part vertex,
must have xor \(03\).  If their entries from \(\{0,2\}\) agree, their xor
is contained in \(\{t_0,t_1\}\), which cannot contain coordinate \(0\).
If those entries differ, their xor contains both \(0\) and \(2\), so it
cannot equal \(03\).  This contradicts the vertex equation at \(y\).
\end{proof}

\begin{remark}
Direct exact checks give full boundary for the adjacent-deletion poles of
the cube, Heawood, M\"obius--Kantor, Pappus, and Desargues graphs, while
\(K_{3,3}\) gives the obstruction above.  These named examples do not
support a theorem for all larger bipartite cubic graphs.  A proof would
need an additional structural hypothesis, and no such hypothesis is
claimed here.
\end{remark}

\section{Replay, trust boundary, and limitations}\label{sec:replay}

The compact source artifacts are:
\begin{verbatim}
scratch/fivecdc-petersen-four-pole-extension-theorem-20260728.md
scratch/verify_petersen_four_pole_extension_theorem_20260728.py
scratch/fivecdc-heawood-full-four-pole-theorem-20260728.md
scratch/verify_heawood_four_pole_full_boundary_20260728.py
scratch/heawood-four-pole-full-boundary-result-20260728.json
scratch/heawood-four-pole-SHA256SUMS-20260728.txt
scratch/verify_k33_four_pole_boundary_20260728.py
\end{verbatim}
Run, from the repository root:
\begin{verbatim}
python3 scratch/verify_petersen_four_pole_extension_theorem_20260728.py
python3 scratch/verify_heawood_four_pole_full_boundary_20260728.py
python3 scratch/verify_k33_four_pole_boundary_20260728.py
sha256sum -c scratch/heawood-four-pole-SHA256SUMS-20260728.txt
\end{verbatim}

The Petersen checker verifies the 13 displayed rows and checks all 100
ordered choices of \(q,r\), all distinct placements encountered by that
ordered enumeration, all 120 coordinate permutations, and every internal
parity equation.  The Heawood checker is separately written and uses only
the Python standard library.  It parses the graph6 record, verifies the
frozen edge table, simplicity, connectivity, bridgelessness, and girth,
checks every displayed row, enumerates all \(10^4\) boundary words, retains
the 640 xor-zero words, and verifies exact coverage by the ten coordinate
orbits.  Its frozen JSON must equal the fresh output.
The \(K_{3,3}\) checker exhausts all \(10^4\) proper-edge labellings of
\(K_{2,2}\), obtains exactly 580 boundary words, and verifies that the
60 missing admissible words are precisely the \(S_5\)-orbit of
\((02,02,03,03)\).

These are positive finite certificates.  They do not rely on a SAT solver,
and no UNSAT proof trace is relevant.  A human can check each row by
placing its edge labels on the displayed edge list and xoring the three
labels at every completed vertex.  The \(K_{3,3}\) negative statement has
both the two-line structural obstruction in Proposition~\ref{prop:k33}
and the complete finite boundary checker.

The limitations are decisive.
\begin{enumerate}
\item Neither insertion theorem supplies a decomposition of every
      bridgeless graph into reducible pieces.
\item A minimum counterexample may contain neither displayed pole.
\item Full boundary for one Heawood pole does not imply full boundary for
      general four-poles, bipartite poles, or high-connectivity poles.
\item The results concern standard Eulerian edge-subsets, not orientable
      cycle double covers.
\item No literature novelty claim is made without specialist review.
\item FiveCDC remains unresolved.
\end{enumerate}

\section*{AI-use statement}

The first-page disclosure is part of this manuscript and must remain in
any derivative circulation.  OpenAI Codex agents materially contributed
to the mathematical search, finite certificates, checking software,
literature triage, and prose.  The human-readable proofs and complete
tables are included so that a reader need not accept an opaque model
output.  A human author should independently verify the tables, rerun the
checkers, complete the prior-art comparison, and accept ordinary scholarly
responsibility before submission.  AI systems must not be listed as
authors.

\bibliographystyle{plain}
\bibliography{references}

\end{document}
